\begin{circuitikz}[
    american,
    voltage dir=RP,
    >=latex,
    %>={Latex[length=5mm, width=2mm]},
    alt={Diagram showing a two-rung ladder logic diagram featuring a Timer Off Delay (TOF) instruction labeled FanTimer and an output coil labeled CoolingFan. On the top rung (Rung 1), the MotorOn contact is open, causing the TOF timer to actively increment its accumulated value (currently 759 out of a 2000 preset). Green power flow highlighting is present on the done (DN) bit of the timer block, indicating that the timer remains done during the timing interval. On the bottom rung (Rung 3), green power flow highlighting spans across the entire rung from the left power rail through the My_Timer.DN contact to energize the CoolingFan output coil.}
    ]

    % Define the number of rungs in the diagram
    \def\numRungs{3}

    % Calculate the length of the vertical power rails dynamically based on number of rungs
    \pgfmathsetmacro{\railLength}{-2 * \numRungs}

    %======================================================================================================
    % Component Grid
    % Spacing = 3x, 2y
    %======================================================================================================
    \foreach \i in {0,...,5} {
        \foreach \j in {0,...,20} { % Change to 0,...,19 if you need 20 individual Y points down to -40
        
        % Calculate spatial positions
        \pgfmathsetmacro{\px}{\i * 3}
        \pgfmathsetmacro{\py}{-1 - (\j * 2)}
        
        % Name coordinate using loop indices: (C-column-row) e.g., (C-0-0), (C-5-10)
        \coordinate (C-\i-\j) at (\px, \py);
        
        % OPTIONAL: Uncomment the 2 lines below to display grid labels while building
        % \fill[red] (C-\i-\j) circle (1.5pt);
        %\node[anchor=south west, scale=0.5, gray] at (C-\i-\j) {(\i,\j)};
        }
    }

    \begin{pgfonlayer}{ft}
        
        \draw[
            BakerBlack,
            very thick,
            name path=L1
        ]
        (0,0)
        -- (0,\railLength);

        \draw[
            BakerBlack,
            very thick,
            name path=L2
        ]
        (15,0)
        -- (15,\railLength);

    \end{pgfonlayer}

    \begin{pgfonlayer}{main}

        %=================================================================
        % Rung 1
        %=================================================================
        \draw[
            BakerBlack
        ]
        (C-0-0)
        to[C, color=BakerBlack, l={{\small MotorOn}}] (C-1-0)
        -- (C-3-0)
        pic (T1) {TON};

        \draw[
            BakerBlack
        ]
        (T1-OUT)
        -- (C-5-0);

        \node[text=BakerBlack, anchor=south] at (T1-TOP) {{\small TOF}};
        \node[text=BakerBlack, anchor=north] at (T1-TOP) {{\small FanTimer}};
        \node[text=BakerBlack, anchor=west] at (T1-PRE) {{\small 2000}};
        \node[text=BakerBlack, anchor=west] at (T1-ACC) {{\small 759}};

        %=================================================================
        % Rung 3 (ton occupies 2 rungs)
        %=================================================================
        \draw[
            BakerBlack
        ]
        (C-0-2)
        to[C, color=BakerBlack, l={{\small My\_Timer.DN}}] (C-1-2)
        -- (C-4-2)
        to[esource, color=BakerBlack, l={{\small CoolingFan}}] (C-5-2);

    \end{pgfonlayer}

    %=================================================================
    % Background layer - highlight any logic that should be true
    % UpnorthGreen, opacity=0.4, line width=10pt
    %=================================================================
    \begin{pgfonlayer}{bg}

        %==================================================================
        % Left vertical rung - leave this alone, it should ALWAYS be green
        % It automatically draws itself to the length defined at the top
        %==================================================================
        \draw[
            UpNorthGreen,
            line width=10pt,
            opacity=0.4
        ]
        (0,0)
        -- (0,\railLength);
        
        %=================================================================
        % Rung 1
        %=================================================================
        \draw[
            UpNorthGreen,
            opacity=0.4,
            line width=10pt,
        ]
        (T1-DN1)
        -- (T1-DN2);

        %=================================================================
        % Rung 3 (ton occupies 2 rungs)
        %=================================================================
        \draw[
            UpNorthGreen,
            opacity=0.4,
            line width=10pt
        ]
        (C-0-2)
        -- (C-5-2);

    \end{pgfonlayer}
    
\end{circuitikz}